\bookchapter{附錄F：摺疊屏選購與使用專題}

\chapterintro{摺疊屏已自“未來科技”變爲“今能買”之成熟品類。然其與直板機於形態、使用習慣、保養方式上差異甚大，值不值得買、何以選、何以用，此篇爲君理清楚。}

\booksection{一、摺疊屏今已發展至何程度？}

先言結論：\textbf{摺疊屏已非“嚐鮮玩具”，然亦非“人人該買”之主流形態}。其介乎小衆極客產品與大衆消費品之間，適特定人羣。

目前主流分兩類：

\begin{itemize}
  \item \textbf{橫向摺疊（書本式）}：合上乃普通手機大小，展開乃接近方形之小平板。代表：華爲Mate X系列、三星Galaxy Z Fold系列、小米MIX Fold系列、vivo X Fold系列。主打"一部手機=手機+平板"。
  \item \textbf{縱向摺疊（翻蓋式）}：合上如粉餅盒，展開乃正常尺寸手機。代表：華爲Pocket系列、三星Galaxy Z Flip系列、OPPO Find N Flip、小米MIX Flip。主打"小巧便攜+外屏速覽"。
\end{itemize}

2024-2026年之關鍵進步：鉸鏈壽命普遍標稱50萬次摺疊（每天開合100次能用13年）、UTG超薄玻璃取代早期塑料膜、摺痕肉眼可見但手感已甚輕微、價格自萬元檔下探至4000-5000元檔。

\begin{tcbinfo}{一句話判斷}
若君今用直板機完全夠用、無“屏幕不夠大”之痛點，即勿買摺疊屏。其乃“錦上添花”之設備，非“雪中送炭”之設備。

\end{tcbinfo}

\booksection{二、何人適買摺疊屏？}

依真實用戶反饋，此幾類人買摺疊屏“真香”概率高：

\begin{enumerate}
  \item \textbf{重度閱讀/文檔處理者}：常於手機上觀PDF、批註文檔、讀長文。展開後接近7.5英寸之內屏，觀A4 PDF不用縮放，體驗質變。
  \item \textbf{經常出差/開會之商務用戶}：不願帶平板又嫌手機屏小。橫向摺疊屏展開觀PPT、籤合同、雙屏對照郵件，比直板機效率高一截。
  \item \textbf{喜歡"懸停態"創作之人}：折至90度放桌上，上半屏取景下半屏按快門（自拍/延時攝影）、上半屏視頻下半屏記筆記、上半屏直播下半屏觀評論——此乃直板機做不到之形態。
  \item \textbf{追求便攜又欲大屏之女性用戶}：縱向摺疊屏合上能塞入小包，展開乃正常手機，外屏還能當化妝鏡、觀通知。
\end{enumerate}

\begin{tcbwarn}{此幾類人勿買}
\begin{itemize}
  \item 手機只用來刷短視頻、聊微信、打電話者——摺疊屏對君毫無增益，反而更重更厚。
  \item 手小或手腕有腱鞘炎者——橫向摺疊屏展開後普遍230g+，單手握持累。
  \item 預算緊張、一部手機要用4年以上者——摺疊屏二手保值率低於直板機，維修成本高。
\end{itemize}

\end{tcbwarn}

\booksection{三、選購要點：觀此6個指標}

勿被“摺疊”二字晃花了眼，選購時重點觀：

\begin{itemize}
  \item \textbf{鉸鏈類型}：水滴鉸鏈（摺痕淺、壽命長，今之主流）優於U型鉸鏈（摺痕深、早期方案）。問銷售或查評測確認爲何種。
  \item \textbf{內屏材質}：必須是UTG超薄玻璃，勿用早期那種塑料膜。UTG手感接近玻璃，塑料膜指甲劃過會有凹陷感且易留痕。
  \item \textbf{摺痕表現}：去線下店親手展開觀。2026年之合格水平乃：亮屏時正面幾觀不出，斜光下能看出一條淺溝，手指劃過有輕微凹陷但不割手。
  \item \textbf{外屏尺寸}：橫向摺疊屏之外屏越大越好（最好6.4英寸以上），合上能當正常手機用；縱向摺疊屏之外屏越大越實用（4英寸以上能回消息、觀導航）。
  \item \textbf{重量與厚度}：橫向摺疊展開後230-260g乃常態，低於230g算優秀；合上厚度12-14mm，低於12mm算薄。縱向摺疊合上後能輕鬆放進襯衫口袋方算合格。
  \item \textbf{售後與保修}：重點觀內屏碎裂之維修價格。有些品牌首年免費換一次內屏，有些自費要2000-4000元。此乃摺疊屏最大之隱性成本。
\end{itemize}

\begin{tcbpractice}{線下體驗清單}
\begin{enumerate}
  \item 親手開合10次，感受鉸鏈阻尼是否順滑、有無異響。
  \item 展開後斜對燈光觀摺痕，再正面亮屏觀之。
  \item 單手持展開之態三十秒，察其重可受與否。
  \item 試用懸停態：折至90度放桌上，試試拍照與觀視頻。
  \item 問清楚內屏單獨維修價格與保修政策。
\end{enumerate}

\end{tcbpractice}

\booksection{四、摺疊屏獨有之使用技巧}

買回來之後，此數個玩法乃直板機體驗不到者：

\booksubsection{分屏與多任務}

橫向摺疊屏既展，多數系統可上下/左右分屏。其常用者：上半屏觀網課視頻，下半屏錄筆記；左微信而右文檔；上地圖導航而下音樂。較直板機之分屏實用遠甚，蓋屏大，每半皆綽綽有餘。

\booksubsection{懸停態（90度摺疊）}

這是摺疊屏的“殺手鐧”用法：

\begin{itemize}
  \item \textbf{自拍/合影}：放桌上折90度，上半屏取景，下半屏按快門，不用尋人幫忙拍。
  \item \textbf{延時攝影/夜景}：同樣放桌上穩定拍攝，省三腳架。
  \item \textbf{視頻通話}：解放雙手，對方觀君之角度亦更自然。
  \item \textbf{看視頻}：放桌上當小電視，飛機/高鐵小桌板上特別實用。
  \item \textbf{外屏預覽+內屏拍攝}：用後置攝像頭拍自己時，外屏能實時顯示畫面，構圖更準。
\end{itemize}

\booksubsection{外屏速覽}

縱向摺疊屏合上後，外屏可以：觀時間/通知/天氣、回微信（語音轉文字或快捷回復）、刷公交卡、觀導航、當化妝鏡（調用前置攝像頭）。不必每次都展開，省電又省事。

\booksubsection{電腦模式}

部分橫向摺疊屏（華爲、三星）展開後可切換“電腦模式”，外接顯示器鍵鼠後像臺簡易PC，臨時處理文檔、回郵件夠用。不是主力生產力工具，但應急很香。

\booksection{五、保養要點：讓摺疊屏活得久一點}

摺疊屏最嬌貴者，內屏與鉸鏈也。此數條習慣能顯著延長其壽命：

\begin{itemize}
  \item \textbf{別在低溫下強行開合}：0℃以下UTG玻璃變脆，開合前先揣懷裏暖幾分鐘。此為官方說明書皆會寫但沒人看之重點。
  \item \textbf{避免指甲/硬物接觸內屏}：內屏雖有保護層但比外屏軟，用指腹操作，別用指甲戳。
  \item \textbf{合上時確認屏內無異物}：一粒沙子夾於中間合上，或徑壓壞內屏。合上前掃一眼。
  \item \textbf{勿貼普通鋼化膜}：內屏只能用官方專用膜或UTG保護膜，普通鋼化膜會影響摺疊與觸控。建議去官方售後貼膜。
  \item \textbf{鉸鏈勿進灰塵液體}：雖多數標稱IPX8防水，然防塵等級普遍較低。沙灘、工地等環境要小心。
  \item \textbf{開合力度均勻}：別單手暴力甩開，雙手自中間勻速展開最穩。
\end{itemize}

\textbf{內屏損壞之高發場景}\par

\begin{itemize}
  \item 合上時夾住耳機線/充電線——直接壓裂內屏。
  \item 放包裏與鑰匙混裝——外屏沒事，內屏被鑰匙硌出裂痕。
  \item 冬天從冷處進暖房立刻開合——溫差導致玻璃應力變化。
  \item 用了一年後覺得"摺痕變深了"去撕原廠膜——撕膜會連內屏一起損壞。
\end{itemize}

\booksection{六、摺疊屏的“坑”和心理預期}

買之前心裏要有這幾本賬：

\begin{itemize}
  \item \textbf{摺痕會一直在}：2026年之技術做不到"完全無痕"，只能做到"亮屏正面幾觀不出"。若汝乃完美主義者，每天盯着摺痕難受，勿買。
  \item \textbf{更重更厚}：橫向摺疊屏展開後普遍230g+，單手握持比直板機累；合上後厚度12-14mm，褲兜裏更鼓。
  \item \textbf{續航普遍一般}：因爲要塞鉸鏈機構，電池空間被壓縮，同價位續航通常不如直板機。
  \item \textbf{保值率低}：一年後二手價通常為新機之50-60\%，比直板機跌得快。
  \item \textbf{維修貴}：內屏總成維修價2000-4000元，鉸鏈維修500-1500元。買之前查清楚保修政策。
  \item \textbf{App適配參差}：主流App已適配大屏，但小衆App可能只為"拉伸放大"，體驗一般。
\end{itemize}

摺疊屏非“更好之手機”，乃“另一種形態之手機”。買它乃因君需其獨特之形態，非因其“更先進”。

—— 關於摺疊屏的選購哲學

\booksection{七、主流品牌何以選（2026年）}

不寫具體型號（會過時），只寫品牌調性，助君快速定位：

\begin{itemize}
  \item \textbf{華爲Mate X系列}：橫向摺疊標杆，鴻蒙生態協同強，外屏比例接近正常手機，合上能當主力機用。價格高然保值率相對最好。
  \item \textbf{三星Galaxy Z Fold/Flip}：摺疊屏鼻祖，鉸鏈調校成熟，Z Flip為縱向摺疊之顏值擔當。國內售後網點較少為短板。
  \item \textbf{小米MIX Fold/Flip}：性價比路線，配置堆料足，價格相對親民。適合想嚐鮮又不想花太多之用戶。
  \item \textbf{vivo X Fold系列}：橫向摺疊裏外屏比例做得最接近正常手機之，合上當直板機用體驗好。
  \item \textbf{OPPO Find N系列}：橫向摺疊裏尺寸偏緊湊之，單手可握持，適合覺得大摺疊太重之用戶。
  \item \textbf{華爲Pocket系列}：縱向摺疊裏外屏最大最實用之，化妝鏡、防曬檢測等女性向功能做得細緻。
\end{itemize}

選購建議：預算充足追求生態選華爲Mate X；欲最成熟鉸鏈選三星Z Fold；性價比首選小米/vivo；縱向摺疊觀華爲Pocket與三星Z Flip。

\begin{tcbinfo}{最後提醒}
摺疊屏技術每年皆於迭代，本附錄寫於2026年。買之前務必查最新評測與用戶反饋，重點關注：鉸鏈壽命實測、摺痕表現、內屏維修政策這三項。線下親手摸永遠比觀評測重要。

\end{tcbinfo}
